\documentclass[12pt,a4paper]{article}

\usepackage[utf8]{inputenc}
\usepackage[T1]{fontenc}
\usepackage{amsmath,amssymb,amsthm}
\usepackage{booktabs}
\usepackage{array}
\usepackage{doi}
\usepackage{hyperref}

\theoremstyle{remark}
\newtheorem*{remark}{Remark}

\newcommand{\Sym}{\mathrm{Sym}}
\newcommand{\SO}{\mathrm{SO}}

\title{The Lorentz-Capacity Sub-Programme\\[4pt]
\normalsize\itshape Projective Origin of the Lorentz Factor in the Cosmochrony Programme}
\author{J\'er\^ome Beau\\[2pt]
\small Independent Researcher\\
\small \texttt{jerome.beau@cosmochrony.org}}
\date{2026}

\begin{document}

  \maketitle

  \begin{abstract}
    The Lorentz-capacity sub-programme answers a specific question in the Cosmochrony
    framework: where does the Lorentz factor $\gamma$ come from, if spacetime geometry is not
    assumed as fundamental?
    The answer is that $\gamma$ need not be introduced as a kinematic postulate, at the price of one named kinematic
    hypothesis.
    It is reconstructed from the bounded Born--Infeld capacity budget of admissible projection,
    then read in the observer-transform structure of special relativity through a supplied Lorentzian co-metric.
    That co-metric has the conditional form of Q11, $\mathrm{diag}(-A_\tau, 2\lambda, 2\lambda, 2\lambda)$, with two
    independent coefficients fixed by no source; it is proportional to the Minkowski form after a rescaling of the
    ordering coordinate, and the normalisation $g^{\mu\nu} = 2\eta^{\mu\nu}$ is a choice of units.
    The link between capacity and kinematics is one named hypothesis, [H-cone]: the light-cone speed of the
    co-metric equals the Born--Infeld bound per ordering step, the coupling datum that Q11 leaves open.

    The kinematic chain, conditional on the supplied co-metric and on [H-cone], with the capacity--metric
    identification left open, is
    \begin{align*}
      &\text{BI budget}
      \;\to\;
      F^\tau = \tfrac{1}{\gamma}
      \;\to\;
      g^{\mu\nu}\propto\eta^{\mu\nu}\ \text{and [H-cone]}
      \;\to\;
      \Lambda(\beta)\\
      &\qquad\to\;
      R(x)^2 = 1-\epsilon(x)
      \;\to\;
      \text{finite-}q\text{ domain.}
    \end{align*}
    Six papers constitute the sub-programme.
    LorCap identifies the projective mobility parameter.
    TempProj constructs the temporal residual map $F^\tau_n =
    \sqrt{1-B_n^2}$.
    LCII extends the construction to the fibrewise non-homogeneous
    regime.
    LC-O1 reconstructs Lorentz transformations and length
    contraction, conditionally on a supplied co-metric and on [H-cone].
    LC-O2-O1 establishes a no-go for the scalar free-fraction
    reduction of the local capacity--metric bridge, which therefore remains open.
    LC-O3 identifies the exact finite-$q$ domain of validity of the
    Carnot parabolic identification.

    \textbf{Interpretive reading.}
    The result is not a new dynamical principle.
    It shows how time dilation, the light cone, Lorentz boosts, and length contraction arise from
    one admissibility-capacity structure once a co-metric of the conditional form of Q11 is supplied and its
    light-cone speed is identified with the Born--Infeld bound.
    It places the structural local lapse and finite-$q$ corrections in the same programme without
    closing the capacity--metric bridge.
  \end{abstract}

  \medskip
  \noindent\textbf{Keywords.}
  Lorentz factor; projective temporal ordering; Born--Infeld capacity; temporal residual map;
  Lorentz transformation; length contraction; capacity--metric bridge; gravitational time
  dilation; finite-$q$ corrections; Cosmochrony.

  \tableofcontents

  \newpage
  \section{Motivation}
    \label{sec:motivation}

    The Lorentz factor is usually introduced as part of the kinematics of special relativity,
    taken as foundational.
    In the Cosmochrony framework this is not satisfactory: the Lorentzian metric is proposed as an
    effective object of admissible non-injective projection.
    It is not yet reconstructed: under [H-L] and [H-lift], Q5b~\cite{Beau2026q5b} obtains a co-metric of rank three,
    and Q11~\cite{Beau2026q11} gives, for a supplied rank-four extension, a form with two independent
    coefficients.
    The structural question is therefore: where does $\gamma = (1-\beta^2)^{-1/2}$ come from?

    The Lorentz-capacity sub-programme answers this by separating two roles that are fused in
    the standard presentation.
    First, a quadratic Born--Infeld capacity budget, a hypothesis expressing the capacity bound
    of~\cite{Beau2026c}, determines the admissible split between temporal and spatial projective capacity.
    Second, a supplied co-metric of the conditional form of Q11 gives the effective Lorentzian interval on which
    observer transformations act as isometries.

    The capacity relation
    \[
      (F^\tau)^2 + \beta^2 = 1
    \]
    is a Euclidean norm constraint in capacity space.
    It is not the Minkowski interval.
    Its role is to determine
    \[
      F^\tau = \sqrt{1-\beta^2} = \frac{1}{\gamma}.
    \]
    The Lorentzian interval is supplied separately, by a co-metric
    \[
      g^{\mu\nu} = \mathrm{diag}(-A_\tau,\, 2\lambda,\, 2\lambda,\, 2\lambda)
    \]
    with $A_\tau$ and $\lambda$ independent, which holds under the hypotheses of Q11 Corollary~6.1 and becomes
    $2\lambda\,\eta^{\mu\nu}$ after the rescaling $\tau' = \tau(2\lambda/A_\tau)^{1/2}$; the normalisation
    $g^{\mu\nu} = 2\eta^{\mu\nu}$ is a choice of units.
    The Lorentz group is therefore not the symmetry group of the capacity sphere.
    It is the isometry group of the effective Lorentzian metric.
    This distinction is the organising principle of the whole sub-programme.
    The two structures meet in one hypothesis, [H-cone]: the capacity load $\beta$, normalised by the Born--Infeld
    bound per depth step, is the velocity of an observer in units of the light-cone speed of the co-metric.
    No source derives it.

    The sub-programme also addresses two downstream questions.
    In the non-homogeneous regime, a local occupancy field $\epsilon(x)$ reduces the local
    residual temporal capacity; the identification of $R(x)^2 = 1-\epsilon(x)$ with the metric
    lapse remains a declared postulate of LCII, LC-O2-O1 having shown that the free-fraction
    reduction does not derive it.
    LC-O3 then establishes the exact finite-$q$ domain in which the Carnot parabolic
    identification of the spatial load is valid.

  \section{The Structural Chain}
    \label{sec:chain}

    Figure~\ref{fig:chain} displays the chain.

    \begin{figure}[h]
      \centering
      \renewcommand{\arraystretch}{1.5}
      \footnotesize
      \begin{tabular}{ccccccccc}
        BI budget & $\to$ & $F^\tau=1/\gamma$ & $\to$ & $g\propto\eta$, [H-cone]
        & $\to$ & $\Lambda(\beta)$ & $\to$ & $R(x)^2=1-\epsilon(x)$\\
        \cite{Beau2026n} & &
        \cite{Beau2026n,Beau2026tp} & &
        \cite{Beau2026q11,Beau2026tp} & &
        \cite{Beau2026lco1} & &
        \cite{Beau2026lco2,Beau2026lco2o1}\\[6pt]
        \multicolumn{7}{c}{} & $\to$ & finite-$q$ \\
        \multicolumn{7}{c}{} & & \cite{Beau2026lco3}
      \end{tabular}
      \caption{The Lorentz-capacity chain.
      Each node cites the paper(s) that state it: the quadratic budget is a hypothesis of LorCap expressing the bound
      [A-cap] of~\cite{Beau2026c}; the co-metric form is that of Q11 Corollary~6.1, supplied under its hypotheses,
      and [H-cone] is the hypothesis of TempProj, which names the coupling datum Q11-O1; the arrow to
      $R(x)^2 = 1-\epsilon(x)$ is a declared postulate.
      The capacity sphere (left) and the effective metric (centre) are independent inputs;
      the Lorentz group acts on the metric, not on the sphere.}
      \label{fig:chain}
    \end{figure}

    Each arrow has a distinct logical status.

    \paragraph{BI budget to temporal residual.}
      The bounded capacity relation fixes the residual temporal capacity once the spatial load
      $\beta$ is known:
      \[
        F^\tau = \sqrt{1-\beta^2}.
      \]
      This is the projective origin of the Lorentz clock factor.
      The quadratic budget is a hypothesis of~\cite{Beau2026n} expressing the capacity bound [A-cap]
      of~\cite{Beau2026c}; its consequence for the mobility is derived in~\cite{Beau2026n}.

    \paragraph{Temporal residual to effective metric.}
      The homogeneous co-metric is supplied: Q5b~\cite{Beau2026q5b} gives rank three under [H-L] and [H-lift],
      and Q11~\cite{Beau2026q11}, Corollary~6.1, gives for a supplied rank-four extension the form
      \[
        g^{\mu\nu} = \mathrm{diag}(-A_\tau,\, 2\lambda,\, 2\lambda,\, 2\lambda)
      \]
      with two independent coefficients.
      Its signature is conditional on those hypotheses; after a rescaling of $\tau$ it is proportional to $\eta$,
      and the normalisation $g^{\mu\nu} = 2\eta^{\mu\nu}$ is a choice of units.
      Neither the global factor nor the coefficient values affect the isometry group.
      This is the point at which Lorentzian geometry enters the chain.
      Reading the load $\beta$ on the light-cone scale of this co-metric is the hypothesis [H-cone].

    \paragraph{Effective metric to Lorentz transformations.}
      Once the metric is available, an inertial observer is a unit timelike splitting of the
      effective interval.
      Under [H-cone], the Lorentz boost $\Lambda(\beta)$ is the unique proper orthochronous isometry mapping one
      splitting to another at capacity load $\beta$ and fixing the transverse directions, established on the tangent
      space in~\cite{Beau2026lco1}.

    \paragraph{Homogeneous capacity to local lapse.}
      The non-homogeneous extension introduces a local occupancy $\epsilon(x)$ and a residual
      capacity radius $R(x)^2 = 1-\epsilon(x)$, defined in~\cite{Beau2026lco2}.
      Paper~\cite{Beau2026lco2o1} shows that the free-fraction law
      $A_\tau = A_\tau^\infty(1-\epsilon)$ does not yield this equivalence: with
      $g^{\tau\tau} = -A_\tau$ it gives $g_{\tau\tau}/g_{\tau\tau}^\infty = 1/(1-\epsilon)$,
      which matches $1-\epsilon$ only at $\epsilon = 0$.
      The bridge therefore remains open, and the equivalence stays a declared postulate
      of~\cite{Beau2026lco2}.

    \paragraph{Finite-$q$ domain.}
      The quadratic identity $B_n^2 + (F_n^\tau)^2 = 1$ is exact when $B_n$ is the exact torus
      load.
      Paper~\cite{Beau2026lco3} identifies the correction $\Delta_q(n) = 4m^2/q^2$ between the
      exact torus count and the Carnot parabolic approximation, and establishes that the approximation is exact for
      $n \leq R := (q-1)/2$ and asymptotically valid for $n = R + o(q)$.

  \section{Paper-by-Paper Synthesis}
  \label{sec:papers}

  \subsection{LorCap: mobility from projective capacity~\texorpdfstring{\cite{Beau2026n}}{}}
    \label{subsec:lorcap}

    LorCap~\cite{Beau2026n} establishes the first link between the admissibility capacity budget
    and relativistic mobility.
    The spatial load is encoded by a projective capacity parameter, and the temporal residual is
    interpreted as the clock-rate factor.
    The central result is the identification of a capacity variable whose Lorentzian reading is
    \[
      F^\tau = \frac{1}{\gamma}.
    \]
    LorCap also establishes why the naive total residual is insufficient: the central coordinate
    of the Heisenberg group contributes only a global phase, so the Lorentz factor must be
    extracted from the temporally projected residual rather than from the raw total capacity.
    The explicit construction of the temporal projection is completed by TempProj.

  \subsection{TempProj: the temporal residual map~\texorpdfstring{\cite{Beau2026tp}}{}}
    \label{subsec:temproj}

    TempProj~\cite{Beau2026tp} constructs the temporal residual map as the unique positive
    function compatible with three inputs: the Lorentzian splitting of a supplied co-metric, the Born--Infeld
    unit capacity sphere, and the exact Weil filling law $\Sigma^{\mathrm{tot}}_n = 1-B_n$, under an identification
    of $B_n$ with the spatial capacity flow stated there.
    The result is
    \[
      P^{Q5b}_\tau\,\Sigma^{\mathrm{tot}}_n = \sqrt{1-B_n^2} = F^\tau_n.
    \]
    Under [H-cone] this gives time dilation, the proper time of an observer at load $\beta$ advancing at the rate
    $F^\tau_n = 1/\gamma$ relative to the rescaled ordering coordinate, and the light cone
    ($B_n = 1 \Rightarrow F^\tau_n = 0$).
    TempProj thereby shows, under its stated hypotheses, that the Lorentz factor can be read as the temporal
    residual of a bounded projective capacity budget.
    The reconstruction of the full transformation between inertial observers is the role of
    LC-O1.

  \subsection{LCII: local fibrewise capacity lapse~\texorpdfstring{\cite{Beau2026lco2}}{}}
    \label{subsec:lcii}

    LCII~\cite{Beau2026lco2} extends the homogeneous construction to the non-homogeneous
    fibrewise regime.
    A local occupancy field $\epsilon(x)$ measures the fraction of projective capacity already
    locked in the fibre over $x$.
    The residual local capacity radius is
    \[
      R(x)^2 := 1-\epsilon(x).
    \]
    The fibrewise temporal map factorises as
    \[
      P^{Q5b}_{\tau,x}\,\Sigma^{\mathrm{tot}}(x,n)
      = R(x)\,\sqrt{f(x,n)\bigl(2-f(x,n)\bigr)},
    \]
    giving a capacity-language account of gravitational time dilation: local occupancy reduces
    the residual temporal capacity.
    LCII introduces two parallel definitions of $R(x)$,
    \[
      R(x)^2 := 1-\epsilon(x), \qquad R(x)^2 := \frac{g_{\tau\tau}(x)}{g_{\tau\tau}^\infty},
    \]
    with a postulated flat-space reference $g_{\tau\tau}^\infty$.
    Their analytic equivalence is the open bridge LC-O2-O1.

  \subsection{LC-O1: Lorentz transformations and length contraction~\texorpdfstring{\cite{Beau2026lco1}}{}}
    \label{subsec:lco1}

    LC-O1~\cite{Beau2026lco1} closes, conditionally, the observer-transform part of the sub-programme.
    The key conceptual clarification is that the Born--Infeld capacity sphere is not a
    Lorentzian structure: it supplies $\beta$ and $1/\gamma$, while the Lorentzian metric is
    a supplied co-metric of the conditional form of Q11, and [H-cone] identifies the load with the velocity of a
    splitting.

    An inertial observer is a future-directed unit timelike vector $u$ satisfying $\eta(u,u) =
  -1$, together with its simultaneity hyperplane $\Sigma_u = \{X \mid \eta(X,u) = 0\}$.
    The unique proper orthochronous isometry of $\eta$ fixing the transverse directions and mapping the rest
    splitting $u_0 = \partial_\tau$ to the moving splitting
    \[
      u_\beta = \gamma(\partial_\tau + \beta\,\partial_x), \qquad \gamma = (1-\beta^2)^{-1/2},
    \]
    is, on the tangent space, the Lorentz boost
    \[
      \tau' = \gamma(\tau - \beta x), \qquad x' = \gamma(x - \beta\tau).
    \]
    Length contraction follows from measuring a rest interval on the moving simultaneity
    hyperplane:
    \[
      \ell = \frac{\ell_0}{\gamma}.
    \]
    Thus LC-O1 shows that projective temporal ordering does not merely produce a scalar clock
    factor.
    Once an effective Lorentzian co-metric is supplied and [H-cone] is assumed, the same capacity load $\beta$
    fixes the complete inertial transformation.

  \subsection{LC-O2-O1: no-go for the free-fraction closure~\texorpdfstring{\cite{Beau2026lco2o1}}{}}
    \label{subsec:lco2o1}

    LC-O2-O1~\cite{Beau2026lco2o1} examines whether a scalar reduction of the local
    admissibility weight closes the non-homogeneous bridge, and establishes that the natural
    candidate does not.
    Under the free-fraction law $A_\tau(x) = A_\tau^\infty\bigl(1-\epsilon(x)\bigr)$, together with
    the principal-symbol identification $g^{\tau\tau} = -A_\tau$ and a diagonal temporal block, the metric
    radius satisfies, for any postulated $A_\tau^\infty > 0$,
    \[
      \frac{g_{\tau\tau}(x)}{g_{\tau\tau}^{\infty}} = \frac{1}{1-\epsilon(x)},
    \]
    which agrees with the capacity definition $R(x)^2 = 1-\epsilon(x)$ only at $\epsilon(x) = 0$.
    What the free-fraction law does support is the co-metric statement
    $g^{\tau\tau}/g^{\tau\tau}_\infty = 1-\epsilon$, the metric inverse of the lapse, which does
    not govern proper-time dilation.
    The result is narrow: a reciprocal scalar law would close the bridge algebraically, so what
    fails is the free-fraction law rather than scalar closure in general.
    LC-O2-O1 is therefore open, and the equivalence of the two definitions of $R(x)$ remains
    the declared postulate of~\cite{Beau2026lco2}.

  \subsection{LC-O3: finite-$q$ domain of the Carnot identification~\texorpdfstring{\cite{Beau2026lco3}}{}}
    \label{subsec:lco3}

    LC-O3~\cite{Beau2026lco3} resolves the saturation-boundary question.
    The quadratic capacity identity $B_n^2 + (F_n^\tau)^2 = 1$ is not corrected: it is exact
    whenever $B_n$ is the exact finite-$q$ torus load.
    The correction concerns only the pre-saturation Carnot parabolic identification
    $B_n^{\mathrm{parab}} = (2n^2+2n+1)/q^2$.

    Setting $R = (q-1)/2$ for odd prime $q$, the exact torus load satisfies
    \[
      B_n^{\mathrm{torus}} = B_n^{\mathrm{parab}} \quad (n \leq R),
      \qquad
      B_n^{\mathrm{torus}} = B_n^{\mathrm{parab}} - \frac{4m^2}{q^2} \quad (n = R+m).
    \]
    The parabolic identification is asymptotically exact for $m/q \to 0$ and has a finite
    limit $\Delta_q \to 4(\alpha - 1/2)^2$ in the macroscopic folded regime $n/q \to \alpha >
  1/2$.
    Thus LC-O3 establishes the domain in which continuum Carnot estimates are valid, without
    weakening the quadratic capacity law.

  \section{Inputs and Outputs}
  \label{sec:deps}

  \subsection{Upstream inputs}

    The sub-programme depends on four upstream inputs.

    \textit{Born--Infeld capacity bound}~\cite{Beau2026c}.
    The quadratic budget hypothesis of~\cite{Beau2026n}, which expresses the bound [A-cap], gives the unit capacity
    form $(F^\tau)^2 + \beta^2 = 1$.
    Its role is to constrain admissible capacity splittings, not to supply a Lorentzian metric.

    \textit{Projective temporal ordering}~\cite{Beau2026pto}.
    The observable $I(n) = \sigma_{\mathrm{pair}}(0) - \sigma_{\mathrm{pair}}(n)$ supplies a
    non-decreasing temporal proxy.
    It supplies no step and no coefficient; the clock-rate reading of $F^\tau$ is that of TempProj, under [H-cone].

    \textit{Supplied co-metric}~\cite{Beau2026q5b,BeauQ8,Beau2026q11}.
    The geometric branch gives the form, not the co-metric: Q5b a rank-three co-metric under [H-L] and [H-lift],
    Q8 the equality $A_H = A_Z = 2\lambda$ with $\lambda$ free for a supplied invariant form, and Q11 the conditional
    rank-four form $\mathrm{diag}(-A_\tau, 2\lambda, 2\lambda, 2\lambda)$ with $A_\tau$ and $\lambda$ independent.
    This is the point at which Lorentzian geometry enters.
    The Lorentz group acts as the isometry group of this effective interval after a rescaling of $\tau$.

    \textit{Cone--capacity identification [H-cone].}
    The light-cone speed of the co-metric equals the Born--Infeld bound per ordering step.
    This is the coupling datum that Q11 names as missing (open problem Q11-O1); the sub-programme assumes it.

  \subsection{Downstream outputs}

    The sub-programme provides three outputs.

    \textit{Special-relativistic kinematics.}
    Time dilation (proper-time rate $1/\gamma$), the light cone ($B = 1 \Rightarrow F^\tau = 0$),
    Lorentz boosts $\Lambda(\beta)$, and length contraction ($\ell = \ell_0/\gamma$) are all
    reconstructed from projective capacity data, a supplied co-metric and [H-cone].

    \textit{Local gravitational capacity lapse.}
    LCII supplies the parallel structural definitions
    $R(x)^2 = 1-\epsilon(x)$ and $R(x)^2 = g_{\tau\tau}(x)/g_{\tau\tau}^\infty$.
    Their equivalence provides a capacity-side reading of gravitational time dilation in the
    Einstein regime, but its derivation from admissibility data remains open.

    \textit{Finite-$q$ control.}
    The exact correction $\Delta_q(n) = 4m^2/q^2$ controls the domain in which continuum Carnot
    estimates of the spatial load can be trusted.

  \section{Status}
  \label{sec:status}

  Table~\ref{tab:status} summarises the logical status of each paper.

  \begin{table}[h]
    \centering
    \caption{Logical status of the Lorentz-capacity sub-programme.}
    \label{tab:status}
    \smallskip
    \footnotesize
    \begin{tabular}{llp{4.5cm}p{4cm}}
      \toprule
      Label & Paper & Main result & Status \\
      \midrule
      LorCap & \cite{Beau2026n} & Lorentz mobility from projective capacity
      & conditional on the quadratic budget and the convergence lemma; kinematic reading under [H-cone] \\
      TempProj & \cite{Beau2026tp} & $F^\tau_n = \sqrt{1-B_n^2}$, time dilation
      & proved under a supplied co-metric, the identification of $B_n$ and the BI lift; time dilation under
      [H-cone] \\
      LCII & \cite{Beau2026lco2} & Fibrewise lapse $R(x)^2 = 1-\epsilon(x)$
      & structural; capacity--metric bridge open \\
      LC-O1 & \cite{Beau2026lco1} & Lorentz boost; length contraction
      & proved for $\eta$; observer reading under a supplied co-metric and [H-cone] \\
      LC-O2-O1 & \cite{Beau2026lco2o1} & no-go for the scalar free-fraction closure
      & proved negative result, for any postulated $A_\tau^\infty$; bridge open \\
      LC-O3 & \cite{Beau2026lco3} & Finite-$q$ correction $\Delta_q = 4m^2/q^2$
      & proved \\
      \bottomrule
    \end{tabular}
  \end{table}

  The Lorentzian co-metric is supplied in the conditional form of Q11, with two independent coefficients, while
  the Born--Infeld capacity relation comes from the quadratic budget hypothesis.
  The inertial kinematic chain is complete conditionally on the supplied co-metric and on [H-cone], and the local
  capacity--metric bridge remains a structural open problem.

  \begin{remark}[Scope of the open bridge]
LC-O2-O1 proves that the scalar free-fraction law does not derive the local identity
$R(x)^2 = 1-\epsilon(x) = g_{\tau\tau}(x)/g_{\tau\tau}^\infty$.
It does not rule out a different scalar or operator-valued closure.
The dynamical equation $G_{\mu\nu} = 8\pi G_N T_{\mu\nu}$ belongs to the spectral entropy
and gravity branch of the programme, while the relation of $\epsilon(x)$ to specific matter
sectors is a cross-branch question.
  \end{remark}

  \section{Conclusion}
  \label{sec:conclusion}

  The Lorentz-capacity sub-programme gives a conditional projective account of the Lorentz
  factor in Cosmochrony.
  The essential result is the chain
  \begin{align*}
    &\text{BI budget}
    \;\to\;
    F^\tau = \tfrac{1}{\gamma}
    \;\to\;
    g^{\mu\nu}\propto\eta^{\mu\nu}\ \text{and [H-cone]}
    \;\to\;
    \Lambda(\beta)\\
    &\qquad\to\;
    R(x)^2 = 1-\epsilon(x)
    \;\to\;
    \text{finite-}q\text{ domain.}
  \end{align*}

  The Born--Infeld capacity sphere supplies the projective load $\beta$ and the residual clock
  factor $1/\gamma$.
  A supplied co-metric of the conditional form of Q11, whose coefficients no source fixes, gives the Lorentzian
  interval, and [H-cone] places the capacity load on its light-cone scale.
  LC-O1 reconstructs Lorentz boosts and length contraction as isometries of that interval.
  LC-O2-O1 shows that the scalar free-fraction reduction does not close the local
  capacity--metric bridge, which remains open.
  LC-O3 proves that finite-$q$ corrections affect the Carnot parabolic identification, not the
  quadratic capacity law.

  In this sense, under the stated hypotheses, the Lorentz factor is read as the temporal residual of a bounded
  admissible projection rather than postulated, and [H-cone] is the one kinematic input that fixes its velocity
  scale.
  The Lorentz group is not imposed on the capacity sphere: it is the isometry group of the
  supplied effective co-metric.
  The Lorentz-capacity sub-programme is therefore complete, conditionally on that co-metric and on [H-cone], as a
  structural bridge between projective temporal ordering and \emph{inertial} relativistic kinematics.
  Its non-homogeneous extension is not: the identification of the local capacity radius with the
  metric lapse remains the declared postulate of~\cite{Beau2026lco2}, LC-O2-O1 having established
  that the scalar free-fraction reduction does not derive it.

  \paragraph{Interpretive outlook.}
  The robust content is therefore a conditional projective account of inertial relativistic kinematics.
  Reading local capacity depletion as gravitational lapse remains a motivated interpretation,
  not a derived result, until LC-O2-O1 is closed by an independently specified local law.

  \bibliographystyle{plain}
  \bibliography{cosmochrony-bibliography}

\end{document}
